\begin{textsample}{POS Dim 6 – ap – Score 7.00 – ap/ap\_em\_40.txt}  \label{ex:f6_pos_109}
CURSO TÉCNICO Os cursos técnicos têm a carga horária mínima de 800 horas, 1.000 horas e até 1.200 horas.
Na rede, ele pode ser ofertado nas formas concomitante, subsequente e integrado — incluindo os da modalidade da Educação de Jovens e Adultos.
CURSO TÉCNICO 800 H 240h Formação básica para o trabalho ( 4 unidades de eixos \textbf{estruturantes} de 80h básica para qualquer habilitação ou eixo tecnológico ) + formação específica + 1 ou mais eixos \textbf{estruturantes} de áreas de conhecimento ou unidades curriculares/eletivas. 1000 H 240h Formação básica para o trabalho ( 4 unidades de eixos \textbf{estruturantes} de 80h básica para qualquer habilitação ou eixo tecnológico ) + formação específica + 1 ou mais eixos \textbf{estruturantes} de áreas de conhecimento ou unidades curriculares/eletivas. 1200 H 240h Formação básica para o trabalho ( 4 unidades de eixos \textbf{estruturantes} de 80h básica para qualquer habilitação ou eixo tecnológico ) + formação específica A carga horária para o estágio profissional supervisionado, quando estipulado no projeto pedagógico do curso, será adicionado à carga horária mínima especificada.
A forma do curso técnico integrado ao ensino médio é ofertada em uma matrícula única para estudantes que concluíram o ensino fundamental.
Os estudantes da forma concomitante têm matrículas distintas para o curso técnico e o ensino médio.
No desenho dos Cursos Técnicos, os eixos \textbf{estruturantes} devem orientar, também, os componentes curriculares da parte de Preparação Básica para o Trabalho da \textbf{matriz} curricular do curso, identificando o foco pedagógico, competências e habilidades.
Os itinerários foram elaborados para substituir a distribuição do conteúdo das 13 disciplinas tradicionais ensinadas ao longo dos três anos do ciclo.
Após a homologação das Diretrizes Curriculares Nacionais do ensino médio e da Base Nacional Comum Curricular ( BNCC ) do Ensino Médio, os itinerários \textbf{formativos} ficaram divididos em cinco áreas de conhecimento : linguagens e suas \textbf{tecnologias} ; matemática e suas \textbf{tecnologias} ; ciências da natureza e suas \textbf{tecnologias} ; ciências humanas e sociais aplicadas e formação técnica e profissional.
Cada um destes itinerários \textbf{formativos} deve se organizar a partir de quatro eixos \textbf{estruturantes} : Investigação Científica, Processos Criativos, \textbf{Mediação} e Intervenção Sociocultural e \textbf{Empreendedorismo}.

% matched lemmas: empreendedorismo, estruturante, formativo, matriz, mediação, tecnologia
\end{textsample}
